
\section{回报：许多有趣的例子}

\textbf{15.5.1.\ 有趣例子：\trProjectiveTransformation。}

$\operatorname{Pic}\mathbb{P}^n_k=\mathbb{Z}$ 这一事实，有着许多美妙而廉价的推论。

\textbf{15.5.A.（练习：射影空间的自同构）}\quad 证明：射影空间 $\mathbb{P}^n_k$ 的全部自同构（固定 $k$ 的）恰对应于 $k$ 上的 $(n+1)\times(n+1)$ 可逆矩阵模去纯量（这个群又称 $\mathrm{PGL}_{n+1}(k)$）。（提示：设 $\pi\colon\mathbb{P}^n_k\to\mathbb{P}^n_k$ 是一个自同构。证明：存在一个同构 $\pi^*\OO(1)\overset{\sim}{\longrightarrow}\OO(1)$，而它进而诱导出一个同构 $\pi^*\colon\Gamma(\mathbb{P}^n,\OO(1))\to\Gamma(\mathbb{P}^n,\OO(1))$。）

\textbf{15.5.2.}\quad 射影空间的自同构常被称为\textbf{\trProjectiveTransformation}（projective transformations）。由于\mbox{练习~15.5.A}，在其「模去纯量的矩阵」这一化身之中，\trProjectiveTransformation{}也被称为\textbf{射影坐标变换}（projective changes of coordinates）。\mbox{练习~15.5.A} 日后会有用，尤其是在 $n=1$ 的情形；此时这些自同构称为\textbf{\trFractionalLinearTransformation}（fractional linear transformations）。（致专家：\mbox{练习~15.5.A} 为什么没有在任意\trBaseRing{} $A$ 上陈述？在那种情形下，论证会在何处出问题？对哪些环 $A$，结果依然成立？）

\textbf{15.5.B.（练习）}\quad 证明：$\Aut(\mathbb{P}^1_k)$ 在 $k$-值点上严格 $3$-可迁；也就是说，给定两组三元组 $(p_1,p_2,p_3)$ 与 $(q_1,q_2,q_3)$，各组都由 $\mathbb{P}^1$ 的互不相同的 $k$-\trValuedPoint{}组成，则恰有一个 $\mathbb{P}^1$ 的自同构把 $p_i$ 送到 $q_i$（$i=1,2,3$）。

\textbf{15.5.C.（练习）}\quad 在任意\trField{} $k$ 上求解下面这些问题。
（a）找出一个\trFractionalLinearTransformation $f(t)\in\mathrm{PGL}(2)$，使得它在 $\mathrm{PGL}(2)$ 中的阶恰好为 $3$。
（b）证明：$\mathrm{PGL}(2)$ 中任意两个阶为 $3$ 的元素都相互共轭。（可能的提示：利用可迁性。）

\textbf{15.5.D.（练习）}\quad 设 $p_0,\dots,p_{n+1}$ 是 $\mathbb{P}^n_k$ 的 $n+2$ 个互不相同的 $k$-\trValuedPoint，其中任何 $n+1$ 个都不落在一张\trHyperplane{}上。证明：存在唯一的\trProjectiveTransformation，把诸 $p_i$（$0\le i\le n$）送到 $[0,\dots,0,1,0,\dots,0]$（其中 $1$ 位于第 $i$ 个位置），并把 $p_{n+1}$ 送到 $[1,\dots,1]$。

\textbf{15.5.E.（有趣的练习）}\quad 设 $X$ 是一个拟射影 $k$-\trScheme，而 $\pi\colon\mathbb{P}^n_k\to X$ 是任意一个（$k$ 上的）\trMorphism。证明：要么 $\pi$ 的像具有维数 $n$，要么 $\pi$ 把 $\mathbb{P}^n_k$ 收缩为一点。特别地，不存在从射影空间到维数更小的拟射影\trVariety{}的非常值映射。提示：证明可以不妨设 $k$ 是\trAlgClosed{}的，从而特别地是无限的。设 $X\subset\mathbb{P}^N$，以 $\pi^*\OO_{\mathbb{P}^N}(1)\cong\OO_{\mathbb{P}^n_k}(d)$ 定义 $d$。试着证明 $d=0$。为此，证明：若 $m\le n$，则 $\mathbb{P}^n$ 中 $m$ 张非空\trHypersurface{}的交非空。在此，利用如下事实：$\mathbb{P}^n_k$ 中任何非空\trHypersurface{}与任何维数至少为 $1$ 的\trSubscheme{}都有非空的交（\mbox{练习~12.3.D}(a)）。

\textbf{15.5.F.}\ \hard\ （练习：写给读过 \S7.6.4——讲群概形的那一\extrahard{}节——的人）\quad 解释 $\mathrm{GL}_n$ 如何（非平凡地！）作用在 $\mathbb{P}^{n-1}$ 上（在 $\mathbb{Z}$ 上，或者在你选定的域上）。（群概形 $\mathrm{GL}_n$ 定义于\mbox{练习~7.6.N}；群概形的\emph{作用}则早见于\mbox{练习~7.6.S}(a)。）提示：比起我们对射影空间的那套开片老描述，动用\trFunctor{}的语言（\S7.6）、采用我们对射影空间的函子式描述（\S15.2.3），这件事要容易做得多。（到\trGrassmannian{}的推广将在\mbox{练习~16.4.L} 中给出。）

\textbf{15.5.3.}\ \emph{\trhRemark{}。}\quad 在\trAlgClosedField{} $\overline{k}$ 上，$\mathrm{GL}_n$ 可迁地作用在 $\mathbb{P}^{n-1}_{\overline{k}}$ 的诸\trClosedPoint{}上，而点 $[1,0,\cdots,0]$ 的\trStabilizer，由「第一列在第一行以下全为 $0$」的矩阵所组成的子群 $P$ 给出。（旁注：把这个点写成列向量更好，这样 $\mathrm{GL}_n$ 的作用便可按通常的方式解读为矩阵乘法。）这暗示 $\mathbb{P}^{n-1}_{\overline{k}}$ 就是商 $\mathrm{GL}_n/P$。这在很大程度上是对的；但我们首先还得把「群的商」这一观念讲得通。

\textbf{15.5.4.\ 更多有趣例子。}

我们现在真的可以动手计算一些\trPicardGroup{}与\trClassGroup{}了。首先是一条有用的观察：注意，你可以把 $Y$ 上的\trInvertibleSheaf{}限制到任何\trSubscheme{} $X$ 上，而这可以是检验一个\trInvertibleSheaf{}非\trTrivial{}的趁手办法。\trEffectiveCartierDivisor（\S9.5.1）有时也能限制：若你在 $Y$ 上有一个\trEffectiveCartierDivisor，则它限制到 $X$ 上是一个闭子概形，局部由单个方程\trSchemeTheoretic{}地切出。如果你运气好，而这个方程又不在 $X$ 的任何\trAssociatedPoint{}上消失（\S6.6.2），那么你便得到 $X$ 上的一个\trEffectiveCartierDivisor。你可以检验：\trEffectiveCartierDivisor{}的限制正对应于\trInvertibleSheaf{}的限制（意义见\mbox{练习~14.1.I}）。

\textbf{15.5.G.（练习：一个\trTorsion\trPicardGroup）}\quad 设 $Y$ 是 $\mathbb{P}^n_k$ 中一张\trHypersurface，对应于一个不可约的次数 $d$ 多项式。证明：$\operatorname{Pic}(\mathbb{P}^n_k-Y)\cong\mathbb{Z}/(d)$。（致微分几何学家：这与 $\pi_1(\mathbb{P}^n_k-Y)\cong\mathbb{Z}/(d)$ 这一事实有关。）提示：\eqref{eq:excision-class-groups}。

接下来两道练习将探究\mbox{练习~15.5.G} 的若干推论，并为我们兑现\mbox{练习~5.4.N} 所许诺的一些例子。

\textbf{15.5.H.（练习：推广\mbox{练习~5.4.N}）}\quad 沿用同一套记号，设 $d>1$（于是 $\operatorname{Pic}(\mathbb{P}^n-Y)\ne 0$），并设 $H_0,\dots,H_n$ 为 $\mathbb{P}^n$ 上的 $n+1$ 张\trCoordinate{}\trHyperplane。证明：$\mathbb{P}^n\setminus Y$ 是\trAffine{}的，且 $\mathbb{P}^n-Y-H_i$ 是它的一个\trDistinguishedOpen{}子集。证明：诸 $\mathbb{P}^n-Y-H_i$ 构成 $\mathbb{P}^n-Y$ 的一个\trOpenCover。证明：$\operatorname{Pic}(\mathbb{P}^n-Y-H_i)=0$。证明：每个 $\mathbb{P}^n-Y-H_i$ 都是某个\trUFD{}的 $\Spec$，而 $\mathbb{P}^n-Y$ 却不是。于是，「是\trUFD」这一性质不是\trAffineLocal{}性质——它只满足\mbox{仿射通信引理~5.3.2} 两条假设中的一条。

\textbf{15.5.I.（练习）}\quad 沿用上一道练习中的同一套记号，证明：在 $\mathbb{P}^n-Y$ 上，$H_i$（限制到这个开集上）是一个不由单个方程切出的\trEffectiveCartierDivisor。（提示：否则它将给出\trClassGroup{}的一个平凡元素。）

\textbf{15.5.J.（练习）}\quad 证明：$A:=\mathbb{R}[x,y]/(x^2+y^2-1)$ 不是\trUFD，但 $A\otimes_{\mathbb{R}}\mathbb{C}$ 是。提示：\mbox{练习~15.5.H}。（反提示：$(1-y)(1+y)=x^2$ 难道不暗示 $A\otimes_{\mathbb{R}}\mathbb{C}$ \emph{不是}\trUFD{}吗？）

\textbf{15.5.K.（练习：$\mathbb{P}^1\times\mathbb{P}^1$ 的\trPicardGroup）}\quad 考虑
\[
X=\mathbb{P}^1_k\times_k\mathbb{P}^1_k\cong\Proj k[w,x,y,z]/(wz-xy),
\]
一张光滑\trQuadricSurface（见\cref{fig:ruling-quadric-surface} 与\mbox{例~10.6.2}）。按下面的思路证明 $\operatorname{Pic}X\cong\mathbb{Z}\oplus\mathbb{Z}$：证明：若 $L=\{\infty\}\times_k\mathbb{P}^1\subset X$ 而 $M=\mathbb{P}^1\times_k\{\infty\}\subset X$，则 $X-L-M\cong\mathbb{A}^2$。这将给你一个满射 $\mathbb{Z}\oplus\mathbb{Z}\twoheadrightarrow\operatorname{Cl}X$。证明：$\OO(L)$ 限制到 $L$ 上是 $\OO$，限制到 $M$ 上是 $\OO(1)$。证明：$\OO(M)$ 限制到 $M$ 上是 $\OO$，限制到 $L$ 上是 $\OO(1)$。（这道练习要花些时间，但极富启发。）

\textbf{15.5.L.（练习）}\quad 证明：不可约的光滑射影\trSurface（$k$ 上的）可以双有理却不同构。提示：利用\trPicardGroup{}证明 $\mathbb{P}^2$ 不同构于 $\mathbb{P}^1\times\mathbb{P}^1$。（旁注：我们将在\mbox{练习~20.2.D} 中看到，「胀开平面」的\trPicardGroup{}是 $\mathbb{Z}^2$；而在\mbox{练习~20.2.E} 中我们将看到，胀开平面并不同构于 $\mathbb{P}^1\times\mathbb{P}^1$，用的是\trPicardGroup{}中再多一点的信息。）

这与\trCurve{}的情形不同：双有理的不可约光滑射影\trCurve（$k$ 上的）必定同构，我们将在\mbox{定理~16.3.3} 中见到。尽管如此，任何两张\trSurface{}都以一种简单的方式相互关联：若 $X$ 与 $X'$ 都是射影的、\trRegular{}的且双有理的，那么把 $X$ 在精心挑选的点上逐次\trBlowup、$X'$ 也如法炮制，便能使两个结果同构（见 [Ha1, Thm.~V.5.5]；\trBlowup{}将在第 22 章讨论）。

\textbf{15.5.M.（练习：\trCone{}的\trPicardGroup）}\quad 设 $X=\Spec k[x,y,z]/(xy-z^2)$，是一个\trCone，其中 $\operatorname{char}k\ne 2$。（特征假设对结果并非必要；之所以把它写进来，是为了让你能用\mbox{练习~5.4.H} 证明 $X$ 的\trNormality。）证明：$\operatorname{Pic}X=0$，而 $\operatorname{Cl}X\cong\mathbb{Z}/2$。提示：证明 $Z=\{x=z=0\}$（「一条直线上的仿射锥」）的类生成 $\operatorname{Cl}X$，办法是证明它的补集 $\mathrm{D}(x)$ 同构于 $\mathbb{A}^2_k$ 的一个\trOpenSubset。证明 $2[Z]=\operatorname{div}(x)$ 从而是\trPrincipal{}的，再证明 $Z$ 不是\trPrincipal{}的（\mbox{练习~15.4.K}）。（注记：你手里的知识已足以证明 $X\setminus\{(0,0,0)\}$ 是\trFactorial{}的。所以，尽管\trClassGroup{}对移去余维数大于 $1$ 的\trLocus{}并不敏感（\S15.4.12），\trPicardGroup{}却并非如此。）

（正规概形上的）一个\trWeilDivisor，若它的某个非零倍数对应于一个\trLineBundle，便称为\textbf{$\mathbb{Q}$-\trCartier{}的}（$\mathbb{Q}$-Cartier）。（过了下一道练习，我们就不再使用这个术语了。）\mbox{练习~15.5.M} 给出了一个\trWeilDivisor{}的例子，它不对应于任何\trLineBundle，却依然是 $\mathbb{Q}$-\trCartier{}的。现在我们就给出一个\emph{不}是 $\mathbb{Q}$-\trCartier{}的\trWeilDivisor{}的例子。

\textbf{15.5.N.（练习：非 $\mathbb{Q}$-\trCartier{}除子）}\quad 在光滑\trQuadricSurface{} $X=\Spec k[w,x,y,z]/(wz-xy)$ 之上的\trCone{}上，设 $Z$ 是由 $w=x=0$ 切出的\trWeilDivisor{}。\mbox{练习~13.1.D} 已经表明：$Z$ 不能由单个方程\trSchemeTheoretic{}地切出。证明更多的内容：若 $n\ne 0$，则 $n[Z]$ 不是\trLocallyPrincipal{}的。从而证明：$Z$ 甚至不能局部地由单个方程按\emph{\trSetTheoretic}的方式切出。提示：利用\mbox{命题~8.3.4}，证明\trAffineScheme{}上一个\trEffectiveCartierDivisor{}的补集也是\trAffine{}的。于是，若 $Z$ 的某个倍数是\trLocallyPrincipal{}的，那么「$Z$ 的补集中由 $y=z=0$ 切出的那个闭子概形」就会是\trAffine{}的——而\trAffineScheme{}的任何闭子概形都是\trAffine{}的。但这正是概形 $y=z=0$（亦称 $wx$-平面）去掉点 $w=x=0$ 之后的差，而我们已见过它是非\trAffine{}的（\S4.4.1）。

\textbf{15.5.O.}\ \hard\ （练习：为具备足够算术背景的读者而设）\quad 把一个\trNumberField{} $K$ 的\trRingOfIntegers{} $\OO_K$ 的（理想）\trClassGroup——如\mbox{练习~14.1.K} 中所定义——与 $\Spec\OO_K$ 的\trClassGroup——如本节中所定义——等同起来。特别地，你将重新得到\trClassGroup{}的常见描述：素理想的形式和，再模去一个来自\trPrincipal{}分式理想的等价关系。

\textbf{15.5.5.\ 再谈\trClassGroup{}与唯一分解。}

如 \S5.4.5 所述，用以检验一个环是否为\trUFD{}的常用手段寥寥无几。下面这道练习便是其中之一，而且确实有用。例如，它蕴含下面这条经典事实：对\trNumberField{}的\trRingOfIntegers{}而言，\trClassGroup{}正是唯一分解的障碍所在（见\mbox{练习~14.1.K} 与\mbox{命题~15.4.16}）。

\textbf{15.5.P.（练习）}\quad 设 $A$ 是一个\trNoetherian\trIntegralDomain。证明：$A$ 是\trUFD，当且仅当 $A$ \trIntegrallyClosed{}且 $\operatorname{Cl}\Spec A=0$。（一个方向是容易的：我们已经证明过\trUFD{}在其\trFractionField{}内\trIntegrallyClosed。另外，\mbox{引理~12.1.7} 表明\trUFD{}的全部余维数 $1$ \trPrimeIdeal{}都是\trPrincipal{}的，这便蕴含 $\operatorname{Cl}\Spec A=0$。剩下要证的，是若 $A$ \trIntegrallyClosed{}且 $\operatorname{Cl}\Spec A=0$，则全部余维数 $1$ \trPrimeIdeal{}都是\trPrincipal{}的，因为这刻画了\trUFD（\mbox{命题~12.3.7}）。你的论证中或许会用到\mbox{代数\trHartogs{}引理~13.5.19}\footnote{原书此处作（Algebraic Hartogs's Lemma 13.5.19, may arise in your argument），「13.5.19」之后的逗号当系衍文，径删。——译者注}。）这是 \S5.4.5 所许诺的\trUFD{}的第三条重要刻画。

我个人最后最喜欢的一个检验环是否为\trUFD{}的办法（\S5.4.5）是\trNagata{}引理。它也是用处最小的一个。

\textbf{15.5.Q.}\ \extrahard\ （练习：\trNagata{}引理）\quad 设 $A$ 是一个\trNoetherian{}整环，$x\in A$ 是一个元素，使得 $(x)$ 是\trPrimeIdeal，而 $A_x=A[1/x]$ 是\trUFD。那么 $A$ 是\trUFD。（提示：\mbox{练习~15.5.P}。利用短正合列\footnote{原书此处作（$\mathbb{Z}[(x)]$），疑当为（$\mathbb{Z}$）（对照 \eqref{eq:excision-class-groups} 的左端），兹照录原书。——译者注}
\[
\mathbb{Z}[(x)]\longrightarrow\operatorname{Cl}\Spec A\longrightarrow\operatorname{Cl}\Spec A_x\longrightarrow 0
\]
（\eqref{eq:excision-class-groups}）来证明 $\operatorname{Cl}\Spec A=0$。先证明 $A[1/x]$ 是\trIntegrallyClosed{}的，然后按如下方式证明 $A$ 也\trIntegrallyClosed。设 $T^n+a_{n-1}T^{n-1}+\cdots+a_0=0$，其中诸 $a_i\in A$，而 $T\in K(A)$。那么由 $A_x$ 的\trIntegralClosedness，我们有 $T=r/x^m$，其中当 $m>0$ 时 $r\notin(x)$。于是若 $m>0$，我们很快便得到矛盾。）

这就引出了\mbox{注记~13.8.6} 所许诺的一个有趣的代数事实。设 $k$ 是特征不为 $2$ 的\trAlgClosedField。令 $A=k[x_1,\dots,x_n]/(x_1^2+\cdots+x_m^2)$，其中 $m\le n$。当 $m\le 2$ 时，我们会遇到一些特殊行为。（若 $m=0$，我们得到\trAffine{}空间；若 $m=1$，我们得到一个非既约概形；若 $m=2$，我们得到一个可约概形，它是两个\trAffine{}空间的并。）

若 $m\ge 3$，我们已在\mbox{练习~5.4.I}(b) 中验证过 $\Spec A$ 是\trNormal{}的。事实上，若 $m\ge 3$，则 $A$ 是\trUFD，\emph{除非} $m=3$ \emph{或者} $m=4$（\mbox{练习~5.4.L}；另见\mbox{练习~13.1.E}）。对 $m=3$ 的情形：
\[
A=k[x,y,z,w_1,\dots,w_{n-3}]/(x^2+y^2-z^2)
\]
不是\trUFD，因为它带有非平凡的\trClassGroup\footnote{原书此处作（as it is has nonzero class group），「is」当系衍词，径删。——译者注}（其论证本质上与\mbox{练习~15.5.M} 相同）。

在 $4$ 处的失效来自二次曲面的几何：我们已经检验过，在 $\Spec k[w,x,y,z]/(wz-xy)$ 中，存在一个余维数 $1$ 的不可约子集——\trRuling{}中一条直线之上的\trCone——它不是\trPrincipal{}的。

\textbf{15.5.R.（练习：$m\ge 5$ 的情形）}\quad 设 $k$ 是特征不为 $2$ 的\trAlgClosedField。证明：若 $\ell\ge 3$，则
\[
A=k[a,b,x_1,\dots,x_n]/(ab-x_1^2-\cdots-x_\ell^2)
\]
是\trUFD，办法是取 $x=a$ 运用\trNagata{}引理。


